\part{Quantum Field Theory}

\section{Relativistic fields and the scalar action}\label{qft:sec:scalar}
A particle trajectory assigns a position to each time. A classical field
assigns a value to each event. For a real scalar field the value lies in
$\mathbb R$; for a spinor field it lies in a complex representation space.
The transformation law of that space is part of the definition of the
field. Quantization will replace the classical field by an
operator-valued distribution; it does not turn a classical field into a
one-particle wavefunction.

We begin with a real scalar field and derive its dynamics before introducing
spin. Compare Srednicki, Chapters~2--3, pp.~15--30
\cite{MSrednicki2007}. The Lorentz representations and spinor algebra below
connect with Chapters~33--36, pp.~205--231; the trace identities connect
with Chapter~47, pp.~294--297. The order here follows the objects and their
relations rather than the order of exercises.

\subsection{The interval and the mass shell}
Work on flat spacetime $\mathbb R^{1,3}$, with $\hbar=c=1$. Use
\begin{equation}\label{qft:metric}
 x^\mu=(t,\mathbf x),\qquad
 \eta_{\mu\nu}=\operatorname{diag}(-1,1,1,1),\qquad
 x_\mu=\eta_{\mu\nu}x^\nu .
\end{equation}
Greek indices range from $0$ to $3$ and spatial indices from $1$ to $3$.
Srednicki writes this metric as $g_{\mu\nu}$; our $\eta_{\mu\nu}$ has the
same signature. In particular,
\begin{equation}\label{qft:interval}
 x\cdot y=\eta_{\mu\nu}x^\mu y^\nu
 =-x^0y^0+\mathbf x\cdot\mathbf y .
\end{equation}
For a nonzero future-directed on-shell four-momentum $p^\mu=(E_{\mathbf p},\mathbf p)$,
\begin{equation}\label{qft:mass-shell}
 p^2=-m^2,\qquad
 E_{\mathbf p}=\sqrt{\mathbf p^2+m^2}>0 .
\end{equation}
Here $m\geq0$ is the mass in natural units. Before quantization it is the
inverse-length parameter of the classical wave equation; the particle
interpretation follows from its energy spectrum.

A Lorentz transformation preserves Eq.~\eqref{qft:interval}. Substituting
$x'^\mu=\Lambda^\mu{}_{\rho}x^\rho$ and
$y'^\nu=\Lambda^\nu{}_{\sigma}y^\sigma$ gives
\begin{align}
 \eta_{\mu\nu}x'^\mu y'^\nu
 &=\eta_{\mu\nu}\Lambda^\mu{}_{\rho}
       \Lambda^\nu{}_{\sigma}x^\rho y^\sigma,\notag\\
 \eta_{\mu\nu}\Lambda^\mu{}_{\rho}
       \Lambda^\nu{}_{\sigma}&=\eta_{\rho\sigma}.
 \label{qft:lorentz-condition}
\end{align}
For $\Lambda^\mu{}_{\nu}=\delta^\mu{}_{\nu}
+\omega^\mu{}_{\nu}+O(\omega^2)$, the terms linear in $\omega$ are
\begin{equation}\label{qft:lorentz-infinitesimal}
 \eta_{\rho\sigma}+\omega_{\sigma\rho}+\omega_{\rho\sigma}
 =\eta_{\rho\sigma},\qquad
 \omega_{\rho\sigma}=-\omega_{\sigma\rho}.
\end{equation}
Thus the infinitesimal transformations have six independent parameters:
three spatial rotations and three boosts.

A scalar field satisfies $\varphi'(x')=\varphi(x)$. The chain rule, with
$x=\Lambda^{-1}x'$, gives
\begin{equation}\label{qft:scalar-gradient}
 \frac{\partial\varphi'}{\partial x'^\mu}(x')
 =(\Lambda^{-1})^\nu{}_{\mu}
   \frac{\partial\varphi}{\partial x^\nu}(x).
\end{equation}
The gradient is a covector. Contracting two such gradients with
$\eta^{\mu\nu}$ produces a scalar by
Eq.~\eqref{qft:lorentz-condition}.

\subsection{Stationary action and the field equation}
For a smooth real field, take the free action
\begin{equation}\label{qft:scalar-action}
 S[\varphi]=\int d^4x\,\mathcal L,\qquad
 \mathcal L=-\frac12\eta^{\mu\nu}
  \frac{\partial\varphi}{\partial x^\mu}
  \frac{\partial\varphi}{\partial x^\nu}
  -\frac12m^2\varphi^2 .
\end{equation}
This is a choice of a local, quadratic, Lorentz-invariant dynamics. Lorentz
invariance alone would also permit interactions such as $\varphi^4$; the
word \emph{free} specifies their absence. With dimensionless action,
$d^4x$ has mass dimension $-4$, so
\begin{equation}\label{qft:scalar-dimensions}
 [\mathcal L]=4,\qquad [\varphi]=1,\qquad [m]=1,
\end{equation}
where brackets denote mass dimension, not a commutator.

Let $f$ be a smooth compactly supported variation and set
$\varphi_\epsilon=\varphi+\epsilon f$. Differentiating at $\epsilon=0$
and using the symmetry of $\eta^{\mu\nu}$,
\begin{align}
 \left.\frac{dS[\varphi_\epsilon]}{d\epsilon}\right|_{\epsilon=0}
 &=\int d^4x\left[
 -\frac12\eta^{\mu\nu}\left(
 \frac{\partial f}{\partial x^\mu}
 \frac{\partial\varphi}{\partial x^\nu}
 +\frac{\partial\varphi}{\partial x^\mu}
 \frac{\partial f}{\partial x^\nu}\right)-m^2\varphi f\right]\notag\\
 &=\int d^4x\left[
 -\eta^{\mu\nu}\frac{\partial\varphi}{\partial x^\nu}
                \frac{\partial f}{\partial x^\mu}-m^2\varphi f\right]
 \notag\\
 &=\int d^4x\left[
 \eta^{\mu\nu}\frac{\partial^2\varphi}
 {\partial x^\mu\partial x^\nu}-m^2\varphi\right]f.
 \label{qft:scalar-variation}
\end{align}
The last line follows by integration by parts; compact support removes the
boundary integral. Since $f$ is arbitrary, stationarity gives
\begin{equation}\label{qft:klein-gordon}
 \left(-\frac{\partial^2}{\partial t^2}
 +\sum_{i=1}^3\frac{\partial^2}{\partial (x^i)^2}-m^2\right)
 \varphi=0 .
\end{equation}
Equivalently, define
$\Box=\eta^{\mu\nu}\frac{\partial^2}
{\partial x^\mu\partial x^\nu}$, so the equation is
$(\Box-m^2)\varphi=0$. Multiplying by $-1$ gives Srednicki's
$(-\partial^2+m^2)\varphi=0$, with his $\partial^2=\Box$.

For $\varphi(x)=A e^{ip\cdot x}$ the derivatives are
\begin{equation}\label{qft:plane-wave}
 \frac{\partial\varphi}{\partial x^\mu}=ip_\mu\varphi,
 \qquad
 (\Box-m^2)\varphi=(-p^2-m^2)\varphi .
\end{equation}
A nonzero plane wave therefore lies on the mass shell
\eqref{qft:mass-shell}. Its positive-frequency phase is
$p\cdot x=-E_{\mathbf p}t+\mathbf p\cdot\mathbf x$.

\subsection{Canonical momentum and energy}
Separating time and space in Eq.~\eqref{qft:scalar-action},
\begin{equation}\label{qft:scalar-lagrangian-time}
 \mathcal L=\frac12\left(\frac{\partial\varphi}{\partial t}\right)^2
 -\frac12|\boldsymbol\nabla\varphi|^2-\frac12m^2\varphi^2 .
\end{equation}
At fixed $\varphi$ and spatial derivatives, the canonical momentum is
\begin{equation}\label{qft:canonical-momentum}
 \pi=\frac{\partial\mathcal L}
 {\partial(\frac{\partial\varphi}{\partial t})}
 =\frac{\partial\varphi}{\partial t}.
\end{equation}
The Legendre transform gives
\begin{align}
 \mathcal H
 &=\pi\frac{\partial\varphi}{\partial t}-\mathcal L\notag\\
 &=\pi^2-\frac12\pi^2
   +\frac12|\boldsymbol\nabla\varphi|^2+\frac12m^2\varphi^2\notag\\
 &=\frac12\left(\pi^2+|\boldsymbol\nabla\varphi|^2+m^2\varphi^2\right),
 \qquad H=\int d^3x\,\mathcal H .\label{qft:scalar-hamiltonian}
\end{align}
For fields with sufficient spatial decay, variation at fixed time and
spatial integration by parts give
\begin{align}
 \delta H&=\int d^3x\left[
 \pi\,\delta\pi+
 \sum_{i=1}^3\frac{\partial\varphi}{\partial x^i}
 \frac{\partial\delta\varphi}{\partial x^i}+m^2\varphi\,\delta\varphi
 \right]\notag\\
 &=\int d^3x\left[
 \pi\,\delta\pi+
 \left(-\sum_{i=1}^3\frac{\partial^2\varphi}{\partial(x^i)^2}
 +m^2\varphi\right)\delta\varphi\right].
 \label{qft:hamiltonian-variation}
\end{align}
Reading off the functional derivatives, Hamilton's equations are
\begin{align}
 \frac{\partial\varphi}{\partial t}
 &=\frac{\delta H}{\delta\pi}=\pi,\notag\\
 \frac{\partial\pi}{\partial t}
 &=-\frac{\delta H}{\delta\varphi}
 =\sum_{i=1}^3\frac{\partial^2\varphi}{\partial(x^i)^2}-m^2\varphi.
 \label{qft:hamilton-field-equations}
\end{align}
Substituting the first equation into the second recovers
Eq.~\eqref{qft:klein-gordon}. The nonnegative energy in
Eq.~\eqref{qft:scalar-hamiltonian} fixes the kinetic sign for this free
scalar theory.

\section{From field modes to particles}\label{qft:sec:quantization}
\subsection{Each free Fourier mode is an oscillator}
For classical fields of sufficient decay write
\begin{equation}\label{qft:spatial-fourier}
 \varphi(t,\mathbf x)=\int\frac{d^3p}{(2\pi)^3}
 \widetilde\varphi(t,\mathbf p)e^{i\mathbf p\cdot\mathbf x}.
\end{equation}
Reality requires
$\widetilde\varphi(t,-\mathbf p)=\widetilde\varphi(t,\mathbf p)^*$.
The spatial second derivatives contribute $-\mathbf p^2$, so
Eq.~\eqref{qft:klein-gordon} gives
\begin{equation}\label{qft:mode-oscillator}
 \frac{\partial^2\widetilde\varphi}{\partial t^2}
 +(\mathbf p^2+m^2)\widetilde\varphi=0.
\end{equation}
Its frequency is precisely $E_{\mathbf p}$ in
Eq.~\eqref{qft:mass-shell}. This is why oscillator quantization enters a
free field theory.

For the oscillator construction take $m>0$, so every momentum has
positive frequency. If $m=0$, the zero-momentum mode needs separate
treatment. In a periodic box of volume $V$, write its spatially constant
field as $\varphi_0(t)=q_0(t)/\sqrt V$. Its Lagrangian and Hamiltonian are
\begin{equation}\label{qft:massless-zero-mode}
 L_0=\frac12\left(\frac{dq_0}{dt}\right)^2,\qquad
 p_{q_0}=\frac{dq_0}{dt},\qquad H_0=\frac12p_{q_0}^2.
\end{equation}
This is a free particle degree of freedom, with no positive-frequency
oscillator vacuum. It cannot be included by setting $E_{\mathbf 0}=0$
in the ladder-operator formulas.

For one oscillator, with $[q,p_q]=i$ and frequency $\omega>0$, define
\begin{equation}\label{qft:oscillator-ladder}
 a=\sqrt{\frac\omega2}\,q+\frac{i}{\sqrt{2\omega}}p_q,
 \qquad
 a^\dagger=\sqrt{\frac\omega2}\,q-\frac{i}{\sqrt{2\omega}}p_q .
\end{equation}
Direct multiplication gives
\begin{align}
 [a,a^\dagger]
 &=-\frac i2[q,p_q]+\frac i2[p_q,q]=1,\notag\\
 a^\dagger a
 &=\frac\omega2q^2+\frac{1}{2\omega}p_q^2
   +\frac i2[q,p_q]
 =\frac\omega2q^2+\frac{1}{2\omega}p_q^2-\frac12,\notag\\
 \frac12(p_q^2+\omega^2q^2)&=\omega\left(a^\dagger a+\frac12\right).
 \label{qft:oscillator-energy}
\end{align}

\subsection{Equal-time commutators and normalization}
Canonical quantization postulates, at equal time,
\begin{equation}\label{qft:canonical-commutators}
 [\varphi(t,\mathbf x),\pi(t,\mathbf y)]
 =i\delta^3(\mathbf x-\mathbf y),\qquad
 [\varphi(t,\mathbf x),\varphi(t,\mathbf y)]
 =[\pi(t,\mathbf x),\pi(t,\mathbf y)]=0.
\end{equation}
Linear fields and their commutators are interpreted as distributions,
with field operators defined by smearing against suitable test functions.
Smearing alone does not define coincident products such as $\pi(x)^2$;
these require a regulator or a specified composite operator, such as a Wick
product. A field at a point is not an ordinary finite-dimensional matrix.

Choose the oscillator normalization
\begin{equation}\label{qft:mode-commutator}
 [a(\mathbf p),a^\dagger(\mathbf q)]
 =(2\pi)^3\delta^3(\mathbf p-\mathbf q),\qquad
 [a(\mathbf p),a(\mathbf q)]=0.
\end{equation}
The Hermitian free field and its momentum are then
\begin{align}
 \varphi(x)
 &=\int\frac{d^3p}{(2\pi)^3\sqrt{2E_{\mathbf p}}}
 \left[a(\mathbf p)e^{ip\cdot x}
       +a^\dagger(\mathbf p)e^{-ip\cdot x}\right],
 \label{qft:field-modes}\\
 \pi(x)
 &=\int\frac{d^3p}{(2\pi)^3}
 \left[-i\sqrt{\frac{E_{\mathbf p}}2}a(\mathbf p)e^{ip\cdot x}
       +i\sqrt{\frac{E_{\mathbf p}}2}a^\dagger(\mathbf p)e^{-ip\cdot x}\right].
 \label{qft:momentum-modes}
\end{align}
The sign in the second line follows by differentiating
$p\cdot x=-E_{\mathbf p}t+\mathbf p\cdot\mathbf x$.
Using Eq.~\eqref{qft:mode-commutator}, the two nonzero terms at equal time
are
\begin{align}
 [\varphi(t,\mathbf x),\pi(t,\mathbf y)]
 &=\frac i2\int\frac{d^3p}{(2\pi)^3}
 \left[e^{i\mathbf p\cdot(\mathbf x-\mathbf y)}
       +e^{-i\mathbf p\cdot(\mathbf x-\mathbf y)}\right]\notag\\
 &=i\delta^3(\mathbf x-\mathbf y).
 \label{qft:mode-normalization-check}
\end{align}
The other two commutators are
\begin{align}
 [\varphi(t,\mathbf x),\varphi(t,\mathbf y)]
 &=\int\frac{d^3p}{(2\pi)^3\,2E_{\mathbf p}}
 \left[e^{i\mathbf p\cdot(\mathbf x-\mathbf y)}
       -e^{-i\mathbf p\cdot(\mathbf x-\mathbf y)}\right]=0,\notag\\
 [\pi(t,\mathbf x),\pi(t,\mathbf y)]
 &=\int\frac{d^3p}{(2\pi)^3}\frac{E_{\mathbf p}}2
 \left[e^{i\mathbf p\cdot(\mathbf x-\mathbf y)}
       -e^{-i\mathbf p\cdot(\mathbf x-\mathbf y)}\right]=0.
 \label{qft:equal-time-zero}
\end{align}
In each integral, changing $\mathbf p\mapsto-\mathbf p$ in the second
term makes it the negative of the first.

Srednicki also uses a covariant oscillator normalization. Write that
oscillator as $b(\mathbf p)$ here to distinguish it from ours:
\begin{align}
 b(\mathbf p)&=\sqrt{2E_{\mathbf p}}\,a(\mathbf p),\notag\\
 [b(\mathbf p),b^\dagger(\mathbf q)]
 &=(2\pi)^3\,2E_{\mathbf p}\delta^3(\mathbf p-\mathbf q),\notag\\
 \varphi(x)&=\int\frac{d^3p}{(2\pi)^3\,2E_{\mathbf p}}
 \left[b(\mathbf p)e^{ip\cdot x}
       +b^\dagger(\mathbf p)e^{-ip\cdot x}\right].
 \label{qft:covariant-normalization}
\end{align}
The change is a normalization of the modes, not a change of the field.
For a massive shell, the measure follows by integrating out positive energy:
\begin{align}
 &\int\frac{d^4p}{(2\pi)^3}\,	heta(p^0)\delta(p^2+m^2)F(p)
 \notag\\
 &\quad=\int\frac{d^3p}{(2\pi)^3}
 \int_0^\infty dp^0\,
 \frac{\delta(p^0-E_{\mathbf p})+\delta(p^0+E_{\mathbf p})}
 {2E_{\mathbf p}}F(p)
 =\int\frac{d^3p}{(2\pi)^3\,2E_{\mathbf p}}
 F(E_{\mathbf p},\mathbf p).
 \label{qft:invariant-shell-measure}
\end{align}
The four-dimensional measure, the mass-shell constraint and the positive-energy
condition are invariant under proper orthochronous Lorentz transformations.
The massless formula holds away from the zero-momentum tip, with its usual
distributional interpretation.

Let $a(\mathbf p)|0\rangle=0$. Evaluate the quadratic quantum
Hamiltonian first in a finite periodic box with a momentum cutoff; sums
then replace the integrals below. Continue to use $m>0$. Normal ordering
places creation operators to the left of annihilation operators and fixes
the free vacuum energy to zero.

To see every contribution, substitute Eqs.~\eqref{qft:field-modes} and
\eqref{qft:momentum-modes} into Eq.~\eqref{qft:scalar-hamiltonian}.
The time derivatives contribute $(-iE_{\mathbf p})(-iE_{\mathbf q})$
to $aa$ and $(-iE_{\mathbf p})(iE_{\mathbf q})$ to
$aa^\dagger$. The spatial derivatives give
$(ip^i)(iq^i)$ and $(ip^i)(-iq^i)$, respectively. Thus define
\begin{align}
 C_-(\mathbf p,\mathbf q)
 &=-E_{\mathbf p}E_{\mathbf q}-\mathbf p\cdot\mathbf q+m^2,
 \notag\\
 C_+(\mathbf p,\mathbf q)
 &= E_{\mathbf p}E_{\mathbf q}+\mathbf p\cdot\mathbf q+m^2.
 \label{qft:hamiltonian-kernel-coefficients}
\end{align}
The conjugate terms have the same respective coefficients. The full
substitution is
\begin{align}
 H&=\int\frac{d^3p\,d^3q}{(2\pi)^6}
 \int\frac{d^3x}{4\sqrt{E_{\mathbf p}E_{\mathbf q}}}
 \mathcal K(p,q;x),\notag\\
 \mathcal K(p,q;x)
 &=C_-(\mathbf p,\mathbf q)\big[
 a(\mathbf p)a(\mathbf q)e^{i(p+q)\cdot x}\notag\\
 &\hspace{5em}+a^\dagger(\mathbf p)a^\dagger(\mathbf q)e^{-i(p+q)\cdot x}
 \big]\notag\\
 &\quad+C_+(\mathbf p,\mathbf q)\big[
 a(\mathbf p)a^\dagger(\mathbf q)e^{i(p-q)\cdot x}\notag\\
 &\hspace{5em}+a^\dagger(\mathbf p)a(\mathbf q)e^{-i(p-q)\cdot x}
 \big].\label{qft:hamiltonian-full-kernel}
\end{align}
For the equal-sign terms,
\begin{align}
 \int d^3x\,e^{i(\mathbf p+\mathbf q)\cdot\mathbf x}
 &=(2\pi)^3\delta^3(\mathbf p+\mathbf q),\notag\\
 \frac{C_-(\mathbf p,-\mathbf p)}{4E_{\mathbf p}}
 &=\frac{-E_{\mathbf p}^2+\mathbf p^2+m^2}{4E_{\mathbf p}}=0.
 \label{qft:hamiltonian-pair-cancellation}
\end{align}
For the mixed terms,
\begin{align}
 \int d^3x\,e^{i(\mathbf p-\mathbf q)\cdot\mathbf x}
 &=(2\pi)^3\delta^3(\mathbf p-\mathbf q),\notag\\
 \frac{C_+(\mathbf p,\mathbf p)}{4E_{\mathbf p}}
 &=\frac{E_{\mathbf p}^2+\mathbf p^2+m^2}{4E_{\mathbf p}}
 =\frac{E_{\mathbf p}}2.
 \label{qft:hamiltonian-mixed-coefficient}
\end{align}
Hence the Hamiltonian and its normal-ordered form are
\begin{align}
 H&=\frac12\int\frac{d^3p}{(2\pi)^3}E_{\mathbf p}
 \left[a(\mathbf p)a^\dagger(\mathbf p)
       +a^\dagger(\mathbf p)a(\mathbf p)\right],\notag\\
 :H:&=\int\frac{d^3p}{(2\pi)^3}E_{\mathbf p}
 a^\dagger(\mathbf p)a(\mathbf p),\notag\\
 [:H:,a^\dagger(\mathbf q)]
 &=\int\frac{d^3p}{(2\pi)^3}E_{\mathbf p}
 a^\dagger(\mathbf p)(2\pi)^3\delta^3(\mathbf p-\mathbf q)
 =E_{\mathbf q}a^\dagger(\mathbf q).
 \label{qft:particle-energy}
\end{align}
The first line carries the regulated vacuum constant. Normal ordering
here fixes the zero of energy in the free flat-spacetime theory; vacuum
effects in interacting theories and in gravity require their own analysis.
Equation~\eqref{qft:particle-energy} gives the particle interpretation:
$a^\dagger(\mathbf q)|0\rangle$ has energy $E_{\mathbf q}$ above the vacuum.
Sharp-momentum states are delta-normalized; normalizable states are wave
packets obtained by integrating them against momentum profiles.

\section{Pauli matrices, Lorentz vectors and two-component spinors}
\label{qft:sec:spinors}
\subsection{The Pauli product}
Take the orientation $\epsilon_{123}=1$ and define
\begin{equation}\label{qft:pauli-matrices}
 \sigma_1=\begin{pmatrix}0&1\\1&0\end{pmatrix},\quad
 \sigma_2=\begin{pmatrix}0&-i\\i&0\end{pmatrix},\quad
 \sigma_3=\begin{pmatrix}1&0\\0&-1\end{pmatrix}.
\end{equation}
Multiplication gives $\sigma_i^2=I_2$ and
\begin{equation}\label{qft:pauli-cyclic}
 \sigma_1\sigma_2=\begin{pmatrix}i&0\\0&-i\end{pmatrix}=i\sigma_3,
 \qquad \sigma_2\sigma_3=i\sigma_1,\qquad
 \sigma_3\sigma_1=i\sigma_2.
\end{equation}
Each reversed product has the opposite sign. Therefore
\begin{equation}\label{qft:pauli-product}
 \sigma_i\sigma_j=\delta_{ij}I_2+i\epsilon_{ijk}\sigma_k,
 \qquad \operatorname{tr}\sigma_i=0.
\end{equation}
For ordinary commuting vectors $\mathbf u,\mathbf v$, this gives
\begin{align}
 (u^i\sigma_i)(v^j\sigma_j)
 &=u^iv^j\delta_{ij}I_2+i u^iv^j\epsilon_{ijk}\sigma_k\notag\\
 &=(\mathbf u\cdot\mathbf v)I_2
   +i(\mathbf u\times\mathbf v)^k\sigma_k.
 \label{qft:pauli-vector-product}
\end{align}
In particular, $(\mathbf p\cdot\boldsymbol\sigma)^2=\mathbf p^2I_2$.

\subsection{A four-vector as a Hermitian matrix}
Associate to a real four-vector $p$ the Hermitian matrix
\begin{equation}\label{qft:hermitian-vector}
 P=p^0I_2+p^i\sigma_i
 =\begin{pmatrix}p^0+p^3&p^1-ip^2\\p^1+ip^2&p^0-p^3\end{pmatrix}.
\end{equation}
Its determinant is the Lorentz norm with a minus sign:
\begin{align}
 \det P&=(p^0+p^3)(p^0-p^3)-(p^1-ip^2)(p^1+ip^2)\notag\\
 &=(p^0)^2-(p^3)^2-(p^1)^2-(p^2)^2=-p^2.
 \label{qft:hermitian-determinant}
\end{align}
For $M\in SL(2,\mathbb C)$ set $P'=MPM^\dagger$. Then
\begin{equation}\label{qft:sl2-action}
 (P')^\dagger=P',\qquad
 \det P'=\det M\,\det P\,\overline{\det M}=\det P.
\end{equation}
The induced real linear transformation preserves the interval. The matrices
$M$ and $-M$ have the same action on $P$, but opposite actions on a spinor
$\lambda\in\mathbb C^2$, which transforms as $\lambda'=M\lambda$.
For example, rotations about the third axis are represented by
\begin{equation}\label{qft:spinor-full-turn}
 M(\theta)=e^{-i\theta\sigma_3/2}
 =\begin{pmatrix}e^{-i\theta/2}&0\\0&e^{i\theta/2}\end{pmatrix},
 \quad M(2\pi)=-I_2,\quad M(4\pi)=I_2 .
\end{equation}
This relates spinors to the $SU(2)$ rotations already developed in the
quaternion part; the spinor itself is not a spatial vector.

\subsection{The alternating form and spinor brackets}
For undotted indices $a,b\in\{1,2\}$, choose
\begin{equation}\label{qft:spinor-epsilon}
 (\varepsilon_{ab})=\begin{pmatrix}0&1\\-1&0\end{pmatrix},\quad
 (\varepsilon^{ab})=\begin{pmatrix}0&-1\\1&0\end{pmatrix},\quad
 \varepsilon^{ab}\varepsilon_{bc}=\delta^a{}_c.
\end{equation}
These conventions fix lowering and raising:
\begin{equation}\label{qft:spinor-raising}
 \lambda_a=\varepsilon_{ab}\lambda^b,
 \qquad \lambda^a=\varepsilon^{ab}\lambda_b,
 \qquad (\lambda_1,\lambda_2)=(\lambda^2,-\lambda^1).
\end{equation}
For $M=\begin{pmatrix}A&B\\C&D\end{pmatrix}$, the explicit invariance
calculation is
\begin{equation}\label{qft:epsilon-invariance}
 M^T\varepsilon M
 =\begin{pmatrix}A&C\\B&D\end{pmatrix}
  \begin{pmatrix}C&D\\-A&-B\end{pmatrix}
 =(AD-BC)\varepsilon=\varepsilon
 \quad (\det M=1).
\end{equation}
Thus the alternating pairing of two commuting spinors is invariant:
\begin{equation}\label{qft:angle-bracket}
 \langle\lambda\,\chi\rangle
 :=\varepsilon_{ab}\lambda^a\chi^b
 =\lambda^1\chi^2-\lambda^2\chi^1
 =-\langle\chi\,\lambda\rangle.
\end{equation}
The word \emph{commuting} matters. Classical fermion fields in a
Grassmann formulation have anticommuting components; exchanging them
introduces an additional sign. The spinors in the determinant identities
here have commuting complex components.

The dimension of the spinor space gives the Schouten identity. In the
first component,
\begin{align}
 &\langle u\,v\rangle w^1+\langle v\,w\rangle u^1
       +\langle w\,u\rangle v^1\notag\\
 &=u^1v^2w^1-u^2v^1w^1+v^1w^2u^1
   -v^2w^1u^1+w^1u^2v^1-w^2u^1v^1=0.
 \label{qft:schouten-component}
\end{align}
The second component cancels independently:
\begin{align}
 &\langle u\,v\rangle w^2+\langle v\,w\rangle u^2
       +\langle w\,u\rangle v^2\notag\\
 &=u^1v^2w^2-u^2v^1w^2+v^1w^2u^2
   -v^2w^1u^2+w^1u^2v^2-w^2u^1v^2=0.
 \label{qft:schouten-component-two}
\end{align}
Hence, without dividing by any bracket,
\begin{equation}\label{qft:schouten-vector}
 \langle u\,v\rangle w^a+\langle v\,w\rangle u^a
       +\langle w\,u\rangle v^a=0.
\end{equation}
Contracting this equation with $\varepsilon_{ab}z^b$ gives
\begin{equation}\label{qft:schouten-four}
 \langle u\,v\rangle\langle w\,z\rangle
 +\langle v\,w\rangle\langle u\,z\rangle
 +\langle w\,u\rangle\langle v\,z\rangle=0.
\end{equation}
It also holds for linearly dependent spinors, where attempts to solve by
dividing by a bracket would fail.

For a nonzero future-directed null momentum, Eq.~\eqref{qft:hermitian-vector}
has eigenvalues $p^0\pm|\mathbf p|=2p^0,0$. It is therefore a positive
rank-one matrix and can be written
\begin{equation}\label{qft:null-spinor-factorization}
 P=\lambda\lambda^\dagger,\qquad
 \lambda^\dagger\lambda=2p^0.
\end{equation}
Indeed, if $e$ is a unit eigenvector for $2p^0$, take
$\lambda=\sqrt{2p^0}\,e$. The phase replacement
$\lambda\mapsto e^{i\alpha}\lambda$ leaves $P$ unchanged.
For complexified momenta, Hermitian conjugation is no longer the relevant
constraint; two independent spinors replace $\lambda$ and its conjugate.
This is the entry to massless spinor-helicity variables; compare Srednicki,
Chapter~50, pp.~308--313 \cite{MSrednicki2007}.

\section{Clifford algebra and Dirac traces}\label{qft:sec:clifford}
\subsection{Factoring the relativistic wave operator}
The scalar wave equation is second order. Seek a first-order matrix
operator whose square supplies its second derivatives. Define
\begin{equation}\label{qft:sigma-four}
 \sigma^\mu=(I_2,\sigma_1,\sigma_2,\sigma_3),\qquad
 \bar\sigma^\mu=(I_2,-\sigma_1,-\sigma_2,-\sigma_3).
\end{equation}
By Eq.~\eqref{qft:pauli-product}, the time-time, mixed and spatial cases
are respectively
\begin{align}
 \sigma^0\bar\sigma^0+\sigma^0\bar\sigma^0&=2I_2,\notag\\
 \sigma^0\bar\sigma^i+\sigma^i\bar\sigma^0&=-\sigma_i+\sigma_i=0,\notag\\
 \sigma^i\bar\sigma^j+\sigma^j\bar\sigma^i
 &=-(\sigma_i\sigma_j+\sigma_j\sigma_i)=-2\delta_{ij}I_2.
 \label{qft:sigma-metric-cases}
\end{align}
Together these say
\begin{equation}\label{qft:sigma-metric}
 \sigma^\mu\bar\sigma^\nu+\sigma^\nu\bar\sigma^\mu
 =-2\eta^{\mu\nu}I_2.
\end{equation}
Exchanging the bars gives the same relation. Form the block matrices
\begin{equation}\label{qft:gamma-definition}
 \gamma^\mu=\begin{pmatrix}0&\sigma^\mu\\\bar\sigma^\mu&0\end{pmatrix}.
\end{equation}
Block multiplication and Eq.~\eqref{qft:sigma-metric} give the Clifford
relation in Srednicki's convention:
\begin{equation}\label{qft:clifford-relation}
 \gamma^\mu\gamma^\nu+\gamma^\nu\gamma^\mu
 =-2\eta^{\mu\nu}I_4.
\end{equation}
In particular $(\gamma^0)^2=I_4$ and $(\gamma^i)^2=-I_4$.

For a smooth four-component spinor $\psi$, commuting partial derivatives
remove the antisymmetric matrix part. Explicitly,
\begin{align}
 \gamma^\mu\gamma^\nu
 \frac{\partial^2\psi}{\partial x^\mu\partial x^\nu}
 &=\frac12(\gamma^\mu\gamma^\nu+\gamma^\nu\gamma^\mu)
 \frac{\partial^2\psi}{\partial x^\mu\partial x^\nu}
 =-\Box\psi,\notag\\
 \left(i\gamma^\mu\frac{\partial}{\partial x^\mu}+m\right)
 \left(i\gamma^\nu\frac{\partial}{\partial x^\nu}-m\right)\psi
 &=(\Box-m^2)\psi.
 \label{qft:dirac-square}
\end{align}
The cross terms cancel because $m$ and the matrices are constant. A
solution of the free Dirac equation
\begin{equation}\label{qft:dirac-equation}
 \left(i\gamma^\mu\frac{\partial}{\partial x^\mu}-m\right)\psi=0
\end{equation}
therefore has Klein--Gordon components. The converse does not follow:
the first-order equation also relates the spinor components.

\subsection{Ordering products without losing metric signs}
The defining relation can be used in either direction:
\begin{equation}\label{qft:gamma-swap}
 \gamma^\mu\gamma^\nu=-\gamma^\nu\gamma^\mu
 -2\eta^{\mu\nu}I_4.
\end{equation}
The second term is essential even when reordering a long product.
Introduce $h^{\mu\nu}=-\eta^{\mu\nu}$ only when comparing the generic
Clifford convention $\{\gamma^\mu,\gamma^\nu\}=2h^{\mu\nu}I_4$.
For the rest of this derivation we use $\eta$ and
Eq.~\eqref{qft:gamma-swap} directly.

Define $\not p=p_\mu\gamma^\mu$. For ordinary commuting four-vectors,
\begin{align}
 (\not p)^2
 &=\frac12p_\mu p_\nu
  (\gamma^\mu\gamma^\nu+\gamma^\nu\gamma^\mu)
 =-p^2I_4,\notag\\
 \not p\not q+\not q\not p&=-2(p\cdot q)I_4.
 \label{qft:slash-product}
\end{align}
On shell, $(\not p)^2=m^2I_4$. The contraction uses the physical
metric $\eta$, even though the Clifford anticommutator contains $-\eta$.

\subsection{Odd traces and the first even traces}
We now use the four-dimensional complex Dirac representation:
$\operatorname{tr}I_4=4$. Define
\begin{equation}\label{qft:chirality}
 \gamma_5=i\gamma^0\gamma^1\gamma^2\gamma^3
 =\begin{pmatrix}-I_2&0\\0&I_2\end{pmatrix},\qquad
 \gamma_5^2=I_4,\qquad
 \gamma_5\gamma^\mu=-\gamma^\mu\gamma_5.
\end{equation}
For a product $G=\gamma^{\mu_1}\cdots\gamma^{\mu_n}$, cyclicity and
Eq.~\eqref{qft:chirality} give
\begin{equation}\label{qft:odd-trace}
 \operatorname{tr}G
 =\operatorname{tr}(\gamma_5G\gamma_5)
 =(-1)^n\operatorname{tr}(G\gamma_5^2)
 =(-1)^n\operatorname{tr}G.
\end{equation}
Thus odd traces vanish. This argument is specifically four-dimensional;
a continuation in dimension requires additional assumptions, particularly
for $\gamma_5$.

For two factors, trace cyclicity and
Eq.~\eqref{qft:clifford-relation} give
\begin{align}
 2\operatorname{tr}(\gamma^\mu\gamma^\nu)
 &=\operatorname{tr}(\gamma^\mu\gamma^\nu+\gamma^\nu\gamma^\mu)
 =-2\eta^{\mu\nu}\operatorname{tr}I_4,\notag\\
 \operatorname{tr}(\gamma^\mu\gamma^\nu)&=-4\eta^{\mu\nu}.
 \label{qft:trace-two}
\end{align}
For four factors, move the first matrix through the remaining three,
using Eq.~\eqref{qft:gamma-swap} at every step:
\begin{align}
 \gamma^\mu\gamma^\nu\gamma^\rho\gamma^\sigma
 &=-\gamma^\nu\gamma^\mu\gamma^\rho\gamma^\sigma
       -2\eta^{\mu\nu}\gamma^\rho\gamma^\sigma\notag\\
 &=\gamma^\nu\gamma^\rho\gamma^\mu\gamma^\sigma
       +2\eta^{\mu\rho}\gamma^\nu\gamma^\sigma
       -2\eta^{\mu\nu}\gamma^\rho\gamma^\sigma\notag\\
 &=-\gamma^\nu\gamma^\rho\gamma^\sigma\gamma^\mu
       -2\eta^{\mu\sigma}\gamma^\nu\gamma^\rho
       +2\eta^{\mu\rho}\gamma^\nu\gamma^\sigma
       -2\eta^{\mu\nu}\gamma^\rho\gamma^\sigma.
 \label{qft:four-gamma-ordering}
\end{align}
The trace of the first term in the last line is minus the original trace.
Bringing it to the left and substituting Eq.~\eqref{qft:trace-two},
\begin{align}
 2\operatorname{tr}(\gamma^\mu\gamma^\nu\gamma^\rho\gamma^\sigma)
 &=-2\eta^{\mu\sigma}(-4\eta^{\nu\rho})
   +2\eta^{\mu\rho}(-4\eta^{\nu\sigma})
   -2\eta^{\mu\nu}(-4\eta^{\rho\sigma}),\notag\\
 \operatorname{tr}(\gamma^\mu\gamma^\nu\gamma^\rho\gamma^\sigma)
 &=4\left(\eta^{\mu\nu}\eta^{\rho\sigma}
 -\eta^{\mu\rho}\eta^{\nu\sigma}
 +\eta^{\mu\sigma}\eta^{\nu\rho}\right).
 \label{qft:trace-four}
\end{align}
For example, $\mu=\nu=0$ in Eq.~\eqref{qft:trace-two} gives $4$, while
$\mu=\nu=i$ with fixed spatial $i$ gives $-4$, as required by the squares
of the matrices. Contracting Eq.~\eqref{qft:trace-four} with four covectors
produces
\begin{align}
 &\operatorname{tr}(\not p\not q\not r\not s)\notag\\
 &\quad=4\left[(p\cdot q)(r\cdot s)-(p\cdot r)(q\cdot s)
                  +(p\cdot s)(q\cdot r)\right].
 \label{qft:slash-trace-four}
\end{align}
This expression is a Lorentz scalar and needs no particular matrix basis.

\subsection{A contracted Clifford identity}
Lower the matrix index with the same physical metric,
$\gamma_\mu=\eta_{\mu\nu}\gamma^\nu$. Contracting the anticommutator,
\begin{align}
 2\gamma^\mu\gamma_\mu
 &=\eta_{\mu\nu}(\gamma^\mu\gamma^\nu+\gamma^\nu\gamma^\mu)
 =-2\eta_{\mu\nu}\eta^{\mu\nu}I_4=-8I_4,\notag\\
 \gamma^\mu\gamma_\mu&=-4I_4.
 \label{qft:gamma-contraction}
\end{align}
Now move $\gamma^\nu$ past $\gamma_\mu$ and apply
Eq.~\eqref{qft:gamma-contraction}:
\begin{align}
 \gamma^\mu\gamma^\nu\gamma_\mu
 &=\gamma^\mu(-\gamma_\mu\gamma^\nu-2\delta^\nu{}_\mu I_4)\notag\\
 &=-\gamma^\mu\gamma_\mu\gamma^\nu-2\gamma^\nu
 =4\gamma^\nu-2\gamma^\nu=2\gamma^\nu.
 \label{qft:gamma-sandwich}
\end{align}
The sign would change if the contracted index were lowered using $h=-\eta$
instead. A contraction is therefore specified by both the algebraic
relation and the bilinear form used to lower the index.
